\clearpage
\fancyhead[L]{\small LINDENAU. ZACH.}
\hypertarget{native-490-lower}{}
\label{native-490-lower}
\begin{center}
\rule{0.55\linewidth}{0.4pt}\\[2pt]
Monatliche Correspondenz. Band XIV. S. 56. 57. Juli 1806.\\[2pt]
\rule{0.55\linewidth}{0.4pt}
\end{center}

\begin{center}
Über \textsc{La Place}s Ausdruck für Höhenmessungen mit dem Barometer. Aus einem Brief eines \emph{früheren} Rezensenten [von \textsc{Lindenau}.]
\end{center}

\begin{flushright}
--- --- --- 18. Mai 1806.
\end{flushright}

--- --- Bei den Ungereimtheiten fällt mir wieder \textsc{Rohde}s letzte in der M.\,C. besprochene Abhandlung ein. N.\,N. und ich waren darüber unterschiedlicher Meinung. N.\,N. erklärte die ganze Methode auch theoretisch für falsch, was ich nicht tat. Dazu schreibt er [\textsc{Gauss}] mir nun Folgendes:

„Auf einen Zweifel, den Sie in Ihrem vorletzten Brief geäußert haben, muss ich noch einmal zurückkommen. Sie finden meine Gründe dafür, dass ich \textsc{Rohde}s Methode auch theoretisch für falsch erklärt habe, zwar einleuchtend. Sie sagen aber, die Übereinstimmung von \textsc{Rohde}s Ausdruck mit einem bekannten Ausdruck habe Sie zögern lassen, die Methode ohne Einschränkung für falsch zu erklären.
\SourceLeaf{491}
Mir scheint, dass dem nur ein kleines Missverständnis zugrunde liegt. Ich habe nicht die von \textsc{Rohde} verwendete Formel für falsch erklärt, sondern seinen Versuch, mit einer zwar richtigen Formel ein Element zu bestimmen. das sich mit dieser Formel nicht bestimmen lässt.“

Darin stimme ich nun auch mit N. N. überein. Ich hätte das nicht erwähnt, wenn Hauptmann \textsc{Rohde} nicht wieder Anlass dazu gegeben hätte usw.

